% =========================================================================
%  Electromagnetismo I — Tema 3: Conductores en equilibrio electrostático
%  Preámbulo. Tipografía, márgenes y espaciado del informe de referencia
%  (clase article, Latin Modern a 11 pt, márgenes de 2,5 cm), con los
%  entornos formales, la notación y el estilo gráfico de los temas 0, 1 y 2.
%  Paleta: azules sobrios con degradados suaves sobre grises neutros.
% =========================================================================

% --- IDIOMA Y CODIFICACIÓN ------------------------------------------------
\usepackage[utf8]{inputenc}
\usepackage[T1]{fontenc}
\usepackage[spanish,es-tabla,es-nodecimaldot,es-noshorthands]{babel}

% --- TIPOGRAFÍA: Latin Modern (Computer Modern escalable), como el informe --
\usepackage{lmodern}
\linespread{1.06}
\emergencystretch=2em
\usepackage[final]{microtype}
\setlength{\parskip}{0.25ex plus 0.15ex}

% --- MATEMÁTICAS -------------------------------------------------------------
\usepackage{amsmath, amssymb, mathtools}
\usepackage{bm}
\usepackage{mathrsfs}
\usepackage{esint}

% --- PÁGINA Y COLORES ---------------------------------------------------------
\usepackage[a4paper, margin=2.5cm, headheight=14pt, headsep=18pt, footskip=30pt]{geometry}
\usepackage{graphicx}
\usepackage[table]{xcolor}
% Paleta del Tema 3: azules sobrios (de la tinta al velo) y grises neutros
\definecolor{azulprof}{HTML}{0B2545}    % azul tinta (títulos, enunciados)
\definecolor{azulmedio}{HTML}{1F4E89}   % azul medio (números, enlaces)
\definecolor{azulacero}{HTML}{3A6EA5}   % azul acero (líneas de campo, acentos)
\definecolor{azulcielo}{HTML}{7FA7D3}   % azul cielo (equipotenciales, marcos)
\definecolor{azulclaro}{HTML}{BCD0EA}   % azul claro (marcos finos)
\definecolor{azulvelo}{HTML}{DCE7F5}    % azul velo (rellenos)
\definecolor{azulfondo}{HTML}{F1F5FB}   % fondo tenue de definiciones
\definecolor{grispizarra}{HTML}{5F6B7A} % gris de texto secundario
\definecolor{grismedio}{HTML}{8A94A3}
\definecolor{grisclaro}{HTML}{D7DDE5}   % rejillas
\definecolor{metal}{HTML}{C9D3E0}       % conductores (relleno)
\definecolor{metalosc}{HTML}{8EA0B8}    % conductores (sombra)
% Acentos mínimos para el signo de las cargas (únicos colores no azules)
\definecolor{cargapos}{HTML}{B03A48}    % carga positiva, vectores destacados
\definecolor{carganeg}{HTML}{1F4E89}    % carga negativa
\definecolor{verdeazul}{HTML}{2A7F8E}   % versores y normales
\colorlet{acento}{cargapos}

% --- TIKZ Y PGFPLOTS ----------------------------------------------------------
\usepackage{tikz}
\usepackage{tikz-3dplot}
\usetikzlibrary{babel, arrows.meta, calc, positioning, decorations.pathmorphing,
                decorations.markings, decorations.pathreplacing, patterns, 3d,
                intersections, angles, quotes, fadings, shapes.geometric, shadings,
                backgrounds, fit}
\usepackage{pgfplots}
\pgfplotsset{compat=1.17}
% Externalización: cada figura se compila una vez a PDF en figuras/ext/ y se
% reutiliza mientras su código no cambie (suma MD5). Requiere -shell-escape
% (ya incluido en latexmkrc; Overleaf lo permite por defecto).
\usepgfplotslibrary{external}
\tikzexternalize[prefix=figuras/ext/]
\tikzset{external/up to date check=md5}
\usepgfplotslibrary{colormaps, fillbetween}
\pgfplotsset{colormap={azulpot}{
  rgb255=(245,248,252) rgb255=(220,231,245) rgb255=(157,185,221)
  rgb255=(58,110,165) rgb255=(31,78,137) rgb255=(11,37,69)}}
\pgfplotsset{colormap={azulcampo}{
  rgb255=(11,37,69) rgb255=(31,78,137) rgb255=(58,134,181)
  rgb255=(127,195,217) rgb255=(207,234,241) rgb255=(255,255,255)}}

% --- TABLAS, LISTAS Y CAJAS ---------------------------------------------------
\usepackage{booktabs}
\usepackage{array}
\usepackage{multirow}
\usepackage{tabularx}
\usepackage{needspace}
\usepackage[section]{placeins}
\usepackage{enumitem}
\setlist[itemize]{leftmargin=1.6em, itemsep=2pt, topsep=4pt, label=\textcolor{azulacero}{\small$\blacktriangleright$}}
\setlist[enumerate]{leftmargin=1.9em, itemsep=2pt, topsep=4pt}
\usepackage[most]{tcolorbox}
\renewcommand{\arraystretch}{1.18}

% --- ENLACES Y REFERENCIAS CRUZADAS ----------------------------------------
% amsthm debe cargarse antes que cleveref para que los enunciados que comparten
% contador se citen con su propio nombre («teorema», «lema»...)
\usepackage{amsthm}
\usepackage[colorlinks=true, linkcolor=azulmedio, citecolor=azulmedio,
            urlcolor=azulmedio, plainpages=false, pdfusetitle=false,
            pdftitle={Tema 3. Conductores en equilibrio electrostático},
            pdfauthor={Daniel Martín Guerrero},
            pdfsubject={Electromagnetismo I}]{hyperref}
\usepackage[spanish,nameinlink]{cleveref}

% --- BIBLIOGRAFÍA -----------------------------------------------------------
\usepackage{csquotes}
\usepackage[backend=biber, style=numeric, sorting=none, giveninits=true,
            maxbibnames=6, minbibnames=3, doi=true, url=true,
            isbn=false]{biblatex}
\addbibresource{referencias.bib}
\DeclareFieldFormat[article,techreport,unpublished,book,inproceedings]{title}{#1\isdot}
\DeclareFieldFormat{url}{\url{#1}}
\renewcommand*{\bibfont}{\small}
\setlength{\bibitemsep}{3pt}

% =========================================================================
%        TÍTULOS DE SECCIÓN, ÍNDICE, CABECERAS Y PIES
% =========================================================================
\usepackage{titlesec}
\usepackage{titletoc}
% Filete con degradado bajo los títulos de sección (no se externaliza)
\newcommand{\filetedegradado}[1][\linewidth]{%
  \par\nointerlineskip\vspace{-0.35\baselineskip}%
  {\tikzset{external/export next=false}%
   \begin{tikzpicture}
     \shade[left color=azulmedio, right color=white, middle color=azulcielo!60]
       (0,0) rectangle (#1,0.75pt);
   \end{tikzpicture}}\par}
\titleformat{\section}[block]
  {\normalfont\Large\bfseries\color{azulprof}}
  {\textcolor{azulmedio}{\thesection.}}{0.7em}{}[\filetedegradado]
\titlespacing*{\section}{0pt}{4.2ex plus 1ex minus .4ex}{1.4ex plus .3ex}
\titleformat{\subsection}[block]
  {\normalfont\large\bfseries\color{azulprof}}
  {\textcolor{azulmedio}{\thesubsection.}}{0.65em}{}
\titlespacing*{\subsection}{0pt}{3.2ex plus 1ex minus .4ex}{1.5ex plus .3ex}
\titleformat{\subsubsection}[block]
  {\normalfont\normalsize\bfseries\color{azulprof}}
  {\textcolor{azulmedio}{\thesubsubsection.}}{0.6em}{}
\titlespacing*{\subsubsection}{0pt}{2.6ex plus .8ex minus .3ex}{1.1ex plus .2ex}
\setcounter{secnumdepth}{3}
\setcounter{tocdepth}{2}

% Índice en la línea del informe: secciones en negrita sin puntos,
% subsecciones con línea de puntos y número de página alineado
\titlecontents{section}[1.9em]{\addvspace{4.5pt}\bfseries\color{azulprof}}
  {\contentslabel[\textcolor{azulmedio}{\thecontentslabel.}]{1.9em}}{}
  {\hfill\textcolor{azulprof}{\contentspage}}
\titlecontents{subsection}[4.6em]{\addvspace{1pt}}
  {\contentslabel[\thecontentslabel.]{2.7em}}{}
  {\titlerule*[0.75pc]{\textcolor{grismedio}{.}}\contentspage}

% Cabeceras discretas en versalitas y número de página centrado (como el informe)
\usepackage{fancyhdr}
\pagestyle{fancy}
\fancyhf{}
\fancyhead[L]{\small\textsc{\color{grispizarra}Electromagnetismo I\enspace·\enspace Tema 3}}
\fancyhead[R]{\small\textsc{\color{grispizarra}\nouppercase{\leftmark}}}
\fancyfoot[C]{\small\color{azulprof}\thepage}
\renewcommand{\headrule}{{\color{azulclaro}\hrule height 0.4pt}\vspace{-0.4pt}}
\renewcommand{\sectionmark}[1]{\markboth{#1}{}}
\fancypagestyle{plain}{\fancyhf{}\fancyfoot[C]{\small\color{azulprof}\thepage}%
  \renewcommand{\headrule}{}}

\numberwithin{equation}{section}

% Pies de figura y tabla
\usepackage[font=small, labelfont={sc,color=azulmedio}, labelsep=period,
            margin=0.6cm, justification=justified]{caption}
\usepackage{subcaption}

% =========================================================================
%               NOTACIÓN Y MACROS DE USO FRECUENTE (temas 0, 1 y 2)
% =========================================================================
% --- Diferencial, derivadas parciales y números -----------------------------
\newcommand{\dif}{\mathrm{d}}
\newcommand{\pd}[2]{\frac{\partial #1}{\partial #2}}
\newcommand{\pdd}[2]{\frac{\partial^{2} #1}{\partial #2^{2}}}
\newcommand{\R}{\mathbb{R}}
\newcommand{\N}{\mathbb{N}}
\newcommand{\eul}{\mathrm{e}}

% --- Vectores, versores y operador nabla ---------------------------------
\newcommand{\vc}[1]{\vec{#1}}                       % vector con flecha
\newcommand{\uh}[1]{\hat{u}_{#1}}                   % versor û_#1
\newcommand{\ux}{\uh{x}}\newcommand{\uy}{\uh{y}}\newcommand{\uz}{\uh{z}}
\newcommand{\ur}{\uh{r}}\newcommand{\uth}{\uh{\theta}}\newcommand{\uph}{\uh{\varphi}}
\newcommand{\urho}{\uh{\rho}}\newcommand{\un}{\hat{n}}
\newcommand{\nab}{\vec{\nabla}}                     % ∇⃗
\newcommand{\Lap}{\nabla^{2}}                       % laplaciano
\newcommand{\norm}[1]{\lVert #1\rVert}
\newcommand{\abs}[1]{\lvert #1\rvert}
\newcommand{\rp}{\vc{r}\,'}                         % r⃗′ punto fuente
\newcommand{\dl}{\dif\vc{\ell}}                     % elemento de línea dirigido
\newcommand{\dS}{\dif\vc{S}}                        % elemento de área dirigido
\newcommand{\dsc}{\dif S}                           % elemento de área escalar
\newcommand{\dV}{\dif V}                            % elemento de volumen
\newcommand{\defeq}{\mathrel{\overset{\mathrm{def}}{=}}}
\newcommand{\bordeS}{\partial S}
\newcommand{\bordeV}{\partial V}
\DeclareMathOperator{\sen}{sen}
\DeclareMathOperator{\sgn}{sgn}
\DeclareMathOperator{\adj}{adj}
\renewcommand{\sin}{\sen}                            % «sen» en todo el texto
\newcommand{\delt}[2]{\delta_{#1#2}}                 % delta de Kronecker
\newcommand{\ord}{\mathcal{O}}                       % orden de Landau

% --- Notación de electrostática (Tema 1) -------------------------------------
\newcommand{\eo}{\varepsilon_{0}}                    % permitividad del vacío
\newcommand{\pot}{\phi}                              % potencial electrostático
\newcommand{\Qenc}{Q_{\mathrm{enc}}}                 % carga encerrada
\newcommand{\Ke}{\frac{1}{4\pi\eo}}                  % constante de Coulomb
\newcommand{\vE}{\vc{E}}
\newcommand{\vF}{\vc{F}}
\newcommand{\dSp}{\dif S\,'}                        % elemento de superficie fuente
\newcommand{\dVp}{\dif V\,'}                        % elemento de volumen fuente
\newcommand{\PhiE}{\Phi_{E}}                         % flujo eléctrico

% --- Notación específica del Tema 3 (conductores) ----------------------------
\newcommand{\Vc}{\mathcal{V}}                        % región ocupada por el conductor
\newcommand{\Sc}{\mathcal{S}}                        % superficie del conductor
\newcommand{\Erest}{\vE_{\mathrm{resto}}}            % campo del resto de la carga
\newcommand{\Edisc}{\vE_{\mathrm{disco}}}            % campo del elemento ΔS
\newcommand{\Pel}{\mathcal{P}}                       % presión electrostática
\newcommand{\Qint}{Q_{\mathrm{int}}}                 % carga de la cara interior
\newcommand{\Qext}{Q_{\mathrm{ext}}}                 % carga de la cara exterior
\newcommand{\Ecav}{\vE_{\mathrm{cav}}}               % campo en la cavidad
\newcommand{\Eext}{\vE_{\mathrm{ext}}}               % campo exterior
\newcommand{\Pm}{\mathsf{P}}                         % matriz de coeficientes de potencial
\newcommand{\Cm}{\mathsf{C}}                         % matriz de coef. de capacidad e influencia
\newcommand{\vQ}{\bm{Q}}                             % vector columna de cargas
\newcommand{\vphi}{\bm{\phi}}                        % vector columna de potenciales
\newcommand{\Erup}{E_{\mathrm{rup}}}                 % campo de ruptura dieléctrica
\newcommand{\DV}{\Delta V}                           % diferencia de potencial
\newcommand{\kc}{k}                                  % k = 1/(4π ε0)
\newcommand{\Tt}{^{\mathsf{T}}}                      % traspuesta

% =========================================================================
%        ENTORNOS FORMALES (enunciados numerados, estilo artículo)
% =========================================================================
% Enunciados en cursiva (resultados generales)
\newtheoremstyle{cienplano}{10pt}{10pt}{\itshape}{}%
  {\bfseries\color{azulprof}}{.}{0.6em}%
  {\thmname{#1}\thmnumber{ #2}\thmnote{\normalfont\upshape\ (#3)}}
% Definiciones, ejemplos y ejercicios (cuerpo en redonda)
\newtheoremstyle{ciendef}{10pt}{10pt}{\normalfont}{}%
  {\bfseries\color{azulprof}}{.}{0.6em}%
  {\thmname{#1}\thmnumber{ #2}\thmnote{\normalfont\upshape\ (#3)}}
% Observaciones y notas (encabezado en cursiva)
\newtheoremstyle{cienobs}{8pt}{8pt}{\normalfont}{}%
  {\itshape\color{azulmedio}}{.}{0.6em}%
  {\thmname{#1}\thmnumber{ #2}\thmnote{\normalfont\upshape\ (#3)}}

\theoremstyle{cienplano}
\newtheorem{proposicion}{Proposición}[section]
\newtheorem{teorema}[proposicion]{Teorema}
\newtheorem{lema}[proposicion]{Lema}
\newtheorem{corolario}[proposicion]{Corolario}
\theoremstyle{ciendef}
\newtheorem{notacion}[proposicion]{Notación}
\newtheorem{ejemplo}[proposicion]{Ejemplo}
\newtheorem{ejercicio}[proposicion]{Ejercicio}
\theoremstyle{cienobs}
\newtheorem{observacion}[proposicion]{Observación}

% Teoremas destacados: filete vertical azul con un velo en degradado
\tcolorboxenvironment{teorema}{enhanced, breakable, shield externalize,
  frame hidden, sharp corners, boxrule=0pt,
  interior style={left color=azulfondo, right color=white},
  borderline west={1.2pt}{0pt}{azulmedio},
  left=8pt, right=4pt, top=0pt, bottom=0pt, before skip=10pt, after skip=10pt}

% Definiciones: caja fina con degradado horizontal y filete lateral.
% Comparten contador con los enunciados (como en los temas anteriores).
\newcounter{definicion}
\renewcommand{\thedefinicion}{\theproposicion}
\newtcolorbox{definicionbox}[1]{%
  enhanced, breakable, sharp corners, shield externalize,
  frame hidden, boxrule=0pt,
  interior style={left color=azulfondo, right color=azulfondo!25!white},
  borderline west={2pt}{0pt}{azulmedio},
  borderline south={0.4pt}{0pt}{azulclaro},
  borderline north={0.4pt}{0pt}{azulclaro},
  left=10pt, right=9pt, top=6pt, bottom=6pt,
  before skip=12pt, after skip=12pt,
  title={Definición~\thedefinicion\ifx\relax#1\relax\else\ {\normalfont(#1)}\fi},
  coltitle=azulprof, fonttitle=\bfseries, detach title, before upper={\textbf{\color{azulprof}\tcbtitle.}\enspace}}
\newenvironment{definicion}[1][]
  {\par\stepcounter{proposicion}\refstepcounter{definicion}\begin{definicionbox}{#1}}
  {\end{definicionbox}}

% Enunciados de los problemas de los seminarios
\newtcolorbox{enunciadobox}[1]{%
  enhanced, breakable, sharp corners, shield externalize,
  frame hidden, boxrule=0pt,
  interior style={top color=azulvelo!55!white, bottom color=azulfondo!40!white},
  borderline west={2pt}{0pt}{azulprof},
  left=10pt, right=9pt, top=6pt, bottom=6pt, before skip=10pt, after skip=10pt,
  fontupper=\itshape,
  title={#1}, coltitle=azulprof, fonttitle=\upshape\bfseries, detach title,
  before upper={{\upshape\bfseries\color{azulprof}\tcbtitle.}\enspace}}
\newcounter{problema}
\renewcommand{\theproblema}{3.\arabic{problema}}
\newenvironment{problema}[1][]
  {\par\refstepcounter{problema}%
   \begin{enunciadobox}{Problema~\theproblema\ifx\relax#1\relax\else\ \normalfont\upshape(#1)\fi}}
  {\end{enunciadobox}}

% Cajas de «receta» y de comprobación (seminarios y guía de métodos)
\newtcolorbox{receta}[1][Receta]{%
  enhanced, breakable, sharp corners, shield externalize,
  colback=white, colframe=azulclaro, boxrule=0.5pt,
  borderline west={2pt}{0pt}{azulacero},
  left=9pt, right=8pt, top=5pt, bottom=5pt, before skip=10pt, after skip=10pt,
  title={#1}, fonttitle=\small\scshape, coltitle=azulprof,
  colbacktitle=azulvelo!70!white, toptitle=1pt, bottomtitle=1pt}

% Ecuaciones destacadas: marco fino con velo en degradado
\tcbset{resaltado/.style={enhanced, shield externalize, sharp corners, nobeforeafter,
  tcbox raise base, colframe=azulcielo, boxrule=0.45pt,
  interior style={left color=azulfondo, right color=white},
  left=4pt, right=4pt, top=2pt, bottom=2pt}}
\newcommand{\resalta}[1]{\tcboxmath[resaltado]{#1}}

% Resolución de ejercicios: misma maquinaria que «proof», con otro rótulo
\renewcommand{\qedsymbol}{\textcolor{azulmedio}{\ensuremath{\blacksquare}}}
\newenvironment{resolucion}[1][Resolución]
  {\begin{proof}[\itshape\color{azulmedio}#1]}{\end{proof}}
% Pasos rotulados dentro de las resoluciones
\newcommand{\paso}[1]{\par\smallskip\noindent\textbf{\color{azulprof}#1}\enspace}

\crefname{proposicion}{proposición}{proposiciones}
\Crefname{proposicion}{Proposición}{Proposiciones}
\crefname{teorema}{teorema}{teoremas}
\Crefname{teorema}{Teorema}{Teoremas}
\crefname{lema}{lema}{lemas}
\Crefname{lema}{Lema}{Lemas}
\crefname{corolario}{corolario}{corolarios}
\Crefname{corolario}{Corolario}{Corolarios}
\crefname{definicion}{definición}{definiciones}
\Crefname{definicion}{Definición}{Definiciones}
\crefname{notacion}{notación}{notaciones}
\Crefname{notacion}{Notación}{Notaciones}
\crefname{ejemplo}{ejemplo}{ejemplos}
\Crefname{ejemplo}{Ejemplo}{Ejemplos}
\crefname{ejercicio}{ejercicio}{ejercicios}
\Crefname{ejercicio}{Ejercicio}{Ejercicios}
\crefname{observacion}{observación}{observaciones}
\Crefname{observacion}{Observación}{Observaciones}
\crefname{problema}{problema}{problemas}
\Crefname{problema}{Problema}{Problemas}
\crefname{section}{sección}{secciones}
\Crefname{section}{Sección}{Secciones}
\crefname{subsection}{sección}{secciones}
\Crefname{subsection}{Sección}{Secciones}
\crefname{subsubsection}{sección}{secciones}
\Crefname{subsubsection}{Sección}{Secciones}
\crefname{table}{tabla}{tablas}
\Crefname{table}{Tabla}{Tablas}
\crefname{figure}{figura}{figuras}
\Crefname{figure}{Figura}{Figuras}
\crefname{equation}{ecuación}{ecuaciones}
\Crefname{equation}{Ecuación}{Ecuaciones}
\crefformat{equation}{ecuación~(#2#1#3)}
\Crefformat{equation}{Ecuación~(#2#1#3)}
\crefrangeformat{equation}{ecuaciones~(#3#1#4)--(#5#2#6)}
\crefmultiformat{equation}{ecuaciones~(#2#1#3)}{ y~(#2#1#3)}{, (#2#1#3)}{ y~(#2#1#3)}

% Referencia a un panel de figura: «figura 5\,(b)», enlazada en bloque
\newcommand{\figpanel}[2]{\hyperref[#1]{figura~\ref*{#1}\,\textup{(#2)}}}
\newcommand{\Figpanel}[2]{\hyperref[#1]{Figura~\ref*{#1}\,\textup{(#2)}}}

% =========================================================================
%        ESTILOS GRÁFICOS COMUNES (TikZ / pgfplots)
% =========================================================================
% =========================================================================
%  Estilos y macros gráficas comunes del Tema 3
%  - \mapcolor{t}: escala cromática azul (t en [0,1]); define el color mapc
%  - Estilos TikZ heredados de los temas 0, 1 y 2 (ejes, vectores, versores…)
%    con la paleta azul del documento
%  - Macros que dibujan los datos numéricos de figuras/datos/*.tex
%    (\LC, \EQ, \FL, \CP, \CN, \SG), generados por scripts/generar_datos.py
% =========================================================================

% --- Escala cromática azul (5 tramos: del velo a la tinta) -------------------
\definecolor{cm0}{HTML}{DCE7F5}
\definecolor{cm1}{HTML}{9DB9DD}
\definecolor{cm2}{HTML}{3A6EA5}
\definecolor{cm3}{HTML}{1F4E89}
\definecolor{cm4}{HTML}{0B2545}
\newcommand{\mapcolor}[1]{%
  \pgfmathsetmacro{\mt}{min(max(#1,0),0.9999)*4}%
  \pgfmathtruncatemacro{\mi}{floor(\mt)}%
  \pgfmathtruncatemacro{\mj}{\mi+1}%
  \pgfmathsetmacro{\mf}{100*(\mt-\mi)}%
  \colorlet{mapc}{cm\mj!\mf!cm\mi}}

% --- Estilos TikZ generales ------------------------------------------------
\tikzset{
  >=Stealth,
  eje/.style={->, line width=0.55pt, grispizarra},
  vect/.style={->, line width=1pt, azulprof},
  vectrojo/.style={->, line width=1pt, cargapos},
  vectverde/.style={->, line width=1pt, verdeazul},
  fino/.style={line width=0.4pt, grispizarra},
  punteado/.style={densely dotted, line width=0.5pt, grispizarra},
  etiq/.style={font=\small, text=black!85},
  cuadro/.style={draw=azulclaro, line width=0.4pt, fill=white, inner sep=2pt},
  % rótulo con fondo blanco semitransparente: evita que el texto se confunda con trazos
  fondo/.style={fill=white, fill opacity=0.85, text opacity=1, inner sep=1.2pt, rounded corners=1pt},
  orienta/.style={decoration={markings, mark=at position 0.55 with {\arrow{Stealth[length=5pt]}}}, postaction={decorate}},
  % conductores: relleno metálico en degradado y contorno en azul tinta
  conductor/.style={draw=azulprof, line width=0.6pt, top color=azulvelo!80!white,
                    bottom color=metal, shading angle=0},
  conductorplano/.style={draw=azulprof, line width=0.6pt, fill=metal!70!white},
  metalrayado/.style={draw=azulprof, line width=0.6pt, fill=azulvelo,
                      postaction={pattern=north east lines, pattern color=azulcielo!70}},
  % superficies gaussianas
  gauss/.style={densely dashed, line width=0.75pt, cargapos!85},
  gaussazul/.style={densely dashed, line width=0.7pt, azulacero},
  % líneas de campo y equipotenciales
  lcampo/.style={azulacero, line width=0.55pt, postaction={decorate},
                 decoration={markings, mark=at position 0.5 with {\arrow{Stealth[length=3.6pt]}}}},
  lcampo2/.style={azulacero, line width=0.55pt, postaction={decorate},
                 decoration={markings, mark=at position 0.3 with {\arrow{Stealth[length=3.6pt]}},
                             mark=at position 0.75 with {\arrow{Stealth[length=3.6pt]}}}},
  lequi/.style={azulcielo, line width=0.45pt, densely dashed},
  % tierra
  tierra/.pic={\draw[line width=0.6pt, azulprof] (0,0) -- (0,-0.18)
                 (-0.18,-0.18) -- (0.18,-0.18) (-0.11,-0.25) -- (0.11,-0.25)
                 (-0.045,-0.32) -- (0.045,-0.32);},
}

% Cargas puntuales y signos de carga superficial
\tikzset{
  qpos/.style={circle, fill=cargapos, draw=cargapos!70!black, line width=0.4pt, inner sep=0pt,
               minimum size=9pt, text=white, font=\scriptsize\bfseries},
  qneg/.style={circle, fill=carganeg, draw=azulprof, line width=0.4pt, inner sep=0pt,
               minimum size=9pt, text=white, font=\scriptsize\bfseries},
  signopos/.style={text=cargapos, font=\scriptsize\bfseries, inner sep=0pt},
  signoneg/.style={text=carganeg, font=\scriptsize\bfseries, inner sep=0pt},
}
\newcommand{\cargaP}[1]{\node[qpos] at (#1) {$+$};}
\newcommand{\cargaN}[1]{\node[qneg] at (#1) {$-$};}

% --- Estilo base de ejes pgfplots (como en los temas anteriores) -----------
\pgfplotsset{
  tema/.style={
    tick label style={font=\scriptsize}, label style={font=\small},
    axis lines=left, axis line style={grispizarra},
    grid=major, grid style={grisclaro, line width=0.35pt},
    every tick/.style={grispizarra},
    legend style={font=\scriptsize, draw=none, fill=white, fill opacity=0.85,
                  text opacity=1},
    legend cell align=left,
    /pgf/number format/use comma,
  },
  marco/.style={
    tick label style={font=\scriptsize}, label style={font=\small},
    axis line style={grispizarra!80, line width=0.5pt},
    every tick/.style={grispizarra}, tick align=inside,
    grid=major, grid style={grisclaro!70, line width=0.3pt},
    /pgf/number format/use comma,
    legend style={font=\scriptsize, draw=none, fill=white, fill opacity=0.85,
                  text opacity=1},
    legend cell align=left,
  },
}

% =========================================================================
%  Macros que interpretan los datos numéricos (figuras/datos/*.tex).
%  Cada figura puede redefinirlas localmente antes de \input.
% =========================================================================
\tikzset{
  numLC/.style={azulmedio, line width=0.45pt},
  numEQ/.style={azulcielo, line width=0.35pt},
  numFL/.style={-{Stealth[length=3.4pt, width=2.8pt]}, azulmedio, line width=0.45pt},
}
\newcommand{\LC}[1]{\draw[numLC] plot coordinates {#1};}
\newcommand{\EQ}[1]{\draw[numEQ] plot coordinates {#1};}
\newcommand{\FL}[3]{\draw[numFL] (#1,#2) -- ++(#3:0.012);}
\newcommand{\CP}[2]{\node[signopos] at (#1,#2) {$+$};}
\newcommand{\CN}[2]{\node[signoneg] at (#1,#2) {$-$};}
\newcommand{\SG}[5]{\mapcolor{#5}\draw[mapc, line width=1.6pt, line cap=round] (#1,#2) -- (#3,#4);}

% Marco fino con marcas (ventana [x0,x1] x [y0,y1]) para los paneles numéricos
\newcommand{\marcopanel}[4]{%
  \draw[azulclaro, line width=0.5pt] (#1,#3) rectangle (#2,#4);}

% =========================================================================
%  Primitivas 3D para tdplot (vista fija: \tdplotsetmaincoords{70}{125}),
%  heredadas de los temas 1 y 2.
% =========================================================================
\def\VA{125}
% \tdCirc[opciones]{radio}{z}   (mitad frontal continua, trasera a trazos)
\newcommand{\tdCirc}[3][]{%
  \draw[#1, domain=\VA-180:\VA, samples=36, variable=\tdt]
    plot ({#2*cos(\tdt)},{#2*sin(\tdt)},{#3});
  \draw[#1, densely dashed, domain=\VA:\VA+180, samples=36, variable=\tdt]
    plot ({#2*cos(\tdt)},{#2*sin(\tdt)},{#3});}
% \tdEsfera{radio}{color}  (esfera sombreada con ecuador)
\newcommand{\tdEsfera}[2]{%
  \begin{scope}[tdplot_screen_coords]
    \shade[ball color=azulvelo, opacity=0.55] (0,0) circle (#1);
    \draw[#2, line width=0.7pt] (0,0) circle (#1);
  \end{scope}
  \tdCirc[#2!60, line width=0.4pt]{#1}{0}}

